\chapter{Phân tích và thiết kế hệ thống}

\section{Yêu cầu hệ thống}
\subsection{Yêu cầu chức năng}
Hệ thống SourceVerify cần đáp ứng các chức năng sau:
\begin{enumerate}
  \item Người dùng có thể tải lên nội dung cần kiểm tra.
  \item Hệ thống xác định loại dữ liệu và kiểm tra tính hợp lệ của tệp.
  \item Với ảnh, hệ thống đọc metadata và ma trận pixel để chạy các method forensic.
  \item Với văn bản và video, hệ thống có cấu trúc mở để bổ sung method tương ứng.
  \item Hệ thống tổng hợp kết quả thành điểm AI, nhãn dự đoán và độ tin cậy.
\end{enumerate}

\subsection{Yêu cầu phi chức năng}
\begin{itemize}
  \item Tính mở rộng: thêm method mới không làm thay đổi toàn bộ kiến trúc.
  \item Tính minh bạch: kết quả phải có giải thích, không chỉ trả về nhãn cuối.
  \item Hiệu năng: ảnh lớn cần được giới hạn kích thước xử lý để tránh quá tải.
  \item Độ tin cậy: kiểm tra định dạng và dữ liệu đầu vào trước khi phân tích.
\end{itemize}

\section{Kiến trúc tổng thể}
Kiến trúc SourceVerify được thiết kế theo hướng đơn giản trước, đúng với nguyên tắc Project I: ưu tiên hiểu rõ bài toán và có sản phẩm minh họa ổn định, sau đó mới mở rộng các thành phần phức tạp hơn.

\begin{figure}[h]
\centering
\begin{tikzpicture}[
  node distance=0.9cm and 1.1cm,
  layer/.style={rectangle,draw=Navy,rounded corners,fill=LightBlue,align=center,minimum width=3.5cm,minimum height=0.9cm},
  arrow/.style={-{Latex[length=2mm]},thick,Navy}
]
\node[layer] (client) {Client Layer\\React UI};
\node[layer,right=of client] (app) {Application Layer\\Next.js Routes};
\node[layer,right=of app] (analysis) {Analysis Layer\\Analyzer Engine};
\node[layer,below=of analysis] (methods) {Method Layer\\Image / Text / Video};
\node[layer,left=of methods] (score) {Scoring Layer\\Weighted Voting};
\node[layer,left=of score] (present) {Presentation Layer\\Result Explainability};
\draw[arrow] (client) -- (app);
\draw[arrow] (app) -- (analysis);
\draw[arrow] (analysis) -- (methods);
\draw[arrow] (methods) -- (score);
\draw[arrow] (score) -- (present);
\draw[arrow] (present) -- (client);
\end{tikzpicture}
\caption{Kiến trúc tổng thể của hệ thống SourceVerify}
\label{fig:system-architecture}
\end{figure}

\section{Phân rã chức năng}
\begin{figure}[h]
\centering
\begin{tikzpicture}[
  root/.style={rectangle,draw=Navy,rounded corners,fill=LightBlue,align=center,minimum width=4.2cm,minimum height=0.8cm},
  child/.style={rectangle,draw=Blue,rounded corners,fill=white,align=center,minimum width=2.9cm,minimum height=0.75cm},
  arrow/.style={-{Latex[length=2mm]},thick,Blue}
]
\node[root] (root) at (0,0) {SourceVerify};
\node[child] (upload) at (-5,-1.5) {Nhập dữ liệu};
\node[child] (pre) at (-2.5,-1.5) {Tiền xử lý};
\node[child] (method) at (0,-1.5) {Phân tích method};
\node[child] (score) at (2.8,-1.5) {Tổng hợp điểm};
\node[child] (ui) at (5.5,-1.5) {Hiển thị kết quả};
\foreach \x in {upload,pre,method,score,ui} \draw[arrow] (root) -- (\x);
\end{tikzpicture}
\caption{Sơ đồ phân rã chức năng của SourceVerify}
\label{fig:function-tree}
\end{figure}

\section{Luồng dữ liệu phân tích ảnh}
Luồng dữ liệu của một lần phân tích ảnh được mô tả trong Hình \ref{fig:data-flow}. Người dùng gửi ảnh lên hệ thống; ảnh được kiểm tra định dạng, đọc metadata, chuyển thành dữ liệu pixel; sau đó các method chạy độc lập và trả về tín hiệu; cuối cùng bộ tổng hợp tạo kết quả.

\begin{figure}[h]
\centering
\begin{tikzpicture}[
  node distance=0.8cm,
  box/.style={rectangle,draw=Navy,rounded corners,fill=LightGray,align=center,minimum width=3.1cm,minimum height=0.75cm},
  data/.style={rectangle,draw=Green,rounded corners,align=center,minimum width=3.1cm,minimum height=0.75cm},
  arrow/.style={-{Latex[length=2mm]},thick,Navy}
]
\node[box] (a) {Upload ảnh};
\node[box,right=of a] (b) {Validate tệp};
\node[data,right=of b] (c) {Metadata + Pixel};
\node[box,below=of c] (d) {Chạy 5 nhóm tín hiệu};
\node[data,left=of d] (e) {Danh sách signal};
\node[box,left=of e] (f) {Weighted scoring};
\node[data,below=of f] (g) {AI score / Verdict};
\node[box,right=of g] (h) {Giải thích kết quả};
\draw[arrow] (a) -- (b);
\draw[arrow] (b) -- (c);
\draw[arrow] (c) -- (d);
\draw[arrow] (d) -- (e);
\draw[arrow] (e) -- (f);
\draw[arrow] (f) -- (g);
\draw[arrow] (g) -- (h);
\end{tikzpicture}
\caption{Luồng dữ liệu phân tích ảnh trong SourceVerify}
\label{fig:data-flow}
\end{figure}

\section{Thiết kế dữ liệu và lựa chọn kiến trúc}
Dữ liệu phân tích được tổ chức theo các cấu trúc logic sau: FileMetadata, AnalysisMethod và AnalysisResult. Thiết kế này giúp kết quả phân tích có cấu trúc thống nhất để bộ tổng hợp có thể dùng lại cho nhiều loại method.

\begin{table}[h]
\centering
\caption{So sánh các lựa chọn kiến trúc cho SourceVerify}
\label{tab:architecture-tradeoff}
\begin{tabularx}{\textwidth}{p{0.22\textwidth}Y Y p{0.14\textwidth}}
\toprule
\textbf{Lựa chọn} & \textbf{Ưu điểm} & \textbf{Nhược điểm} & \textbf{Độ phức tạp} \\
\midrule
Monolithic Next.js + module analyzer & Dễ triển khai, dễ hiểu, phù hợp Project I, ít phụ thuộc hạ tầng & Khó mở rộng khi tải lớn & Thấp \\
Microservice analyzer riêng & Dễ scale từng service, phù hợp sản phẩm lớn & Cần hạ tầng, queue, API nội bộ, monitoring & Cao \\
Mô hình ML end-to-end & Có thể đạt độ chính xác cao nếu có dataset lớn & Cần dữ liệu gán nhãn, training, đánh giá nghiêm ngặt & Cao \\
\bottomrule
\end{tabularx}
\end{table}

Quyết định hiện tại là dùng kiến trúc Next.js kết hợp analyzer module nội bộ. Lý do là nhóm thực hiện cần tập trung vào hiểu bài toán và trình bày method, trong khi microservice hoặc training model sẽ làm tăng độ phức tạp vượt phạm vi Project I.
